\documentclass[11pt]{article}
\usepackage[T1]{fontenc}
\usepackage[margin=1in]{geometry}
\usepackage{amsmath,amssymb,amsthm,mathtools,booktabs,array}
\usepackage[colorlinks=true,linkcolor=blue,urlcolor=blue,citecolor=blue]{hyperref}
\usepackage{microtype}
\hypersetup{pdftitle={Local link replacement with explicit boundary conditions},pdfauthor={Research note and proof audit},pdfsubject={Finite SU(2) Wilson estimates and an obstruction; not a continuum Yang--Mills proof}}
\newtheorem{theorem}{Theorem}[section]
\newtheorem{lemma}[theorem]{Lemma}
\newtheorem{proposition}[theorem]{Proposition}
\newtheorem{corollary}[theorem]{Corollary}
\theoremstyle{definition}
\newtheorem{definition}[theorem]{Definition}
\newtheorem{assumption}[theorem]{Assumption}
\theoremstyle{remark}
\newtheorem{remark}[theorem]{Remark}
\newcommand{\Tr}{\operatorname{Tr}}
\newcommand{\E}{\mathbb E}
\newcommand{\R}{\mathbb R}
\newcommand{\one}{\mathbf1}
\newcommand{\norm}[1]{\left\lVert#1\right\rVert}
\newcommand{\dd}{\,\mathrm d}
\newcommand{\SU}{\mathrm{SU}(2)}
\setlength{\emergencystretch}{3em}
\title{Local link replacement with explicit boundary conditions}
\author{Research note and proof audit}
\date{October 9, 2026}
\begin{document}
\maketitle
\begin{abstract}
We examine the rough-region replacement step suggested by Family 215's
spin-model argument. An unconditional extension that preserves arbitrary
frozen gauge boundary links while removing every rough plaquette is false:
we exhibit an open set of frustrated one-link boundaries. We prove instead
an excess-energy tail estimate for every boundary, and an explicit
measurable gauge-covariant replacement kernel for a sufficient class of
compatible boundaries. For $n$ edited links in four-dimensional $SU(2)$
Wilson theory at inverse coupling $H\ge1$, the latter kernel has density
at most $(24H^{3/2})^n$, makes every affected plaquette curvature at most
$4/\sqrt H$, and gives a conditional energy--entropy bound
$\exp[-kp^2/2+n(48+\log24+\tfrac32\log H)]$ for $k$ plaquettes of
curvature greater than $p/\sqrt H$.
The compatibility condition is explicit, gauge invariant, sufficient but
not necessary. No estimate is supplied for the probability of incompatible
boundaries, for arbitrary renormalized interactions, or for cutoff
comparison. This is not a continuum construction or a mass-gap proof,
and no novelty claim is made.
\end{abstract}
\tableofcontents

\section{Finite lattice and the exact task}
Work on a four-dimensional rectangular periodic lattice whose periods are
at least four. Each positively oriented link $e$ has a variable
$U_e\in\SU$; its reverse has variable $U_e^{-1}$. Each unoriented elementary
plaquette is counted once with a fixed orientation. Set
\begin{equation}
 s(g)=1-\tfrac12\Tr g,\qquad
 E(U)=\sum_p s(U_p),\qquad
 \mu_H(\dd U)=Z_H^{-1}e^{-HE(U)}\prod_e\dd U_e,
 \quad H\ge1,
 \label{eq:model}
\end{equation}
with probability Haar measure on every link. The trace of an $SU(2)$
matrix is real. The nonnegative-action convention here differs from the
unshifted Wilson convention only by an extensive constant.

Fix, \emph{before integrating}, a nonempty set $J$ of editable links, and
write $n=|J|$, $x=U_J$, $z=U_{J^c}$. Let $P(J)$ consist of all plaquettes
that touch at least one link in $J$, and let $q_p=|\partial p\cap J|$.
Define
\[
 E_J(x,z)=\sum_{p\in P(J)}s(U_p).
\]
Each link belongs to six plaquettes, so
\begin{equation}
 |P(J)|\le6n,\qquad \sum_{p\in P(J)}q_p=6n,\qquad
 \sum_{p\in P(J)}q_p^2\le24n.\label{eq:incidence}
\end{equation}
For an open finite complex with all affected plaquettes included, the
corresponding statements are inequalities. Only $E_J$ changes when $x$
is replaced. The conditional law given $z$ is proportional to
$e^{-HE_J(x,z)}\dd x$.

The original spin replacement in \cite{repo-o3} combines exterior
compatibility, a small-density-cost sampling kernel, and a lower bound
on original energy proportional to the edited volume. These are distinct
requirements. In a gauge theory one must edit \emph{links}, not independent
plaquette matrices; this automatically respects all algebraic identities
among plaquettes, but does not guarantee small plaquette curvature.

\section{Geometry and Haar-volume constants}
Use the chordal operator metric
$d(g,h)=\norm{g-h}_{\mathrm{op}}$. It is bi-invariant and obeys
\[
 d(\one,g_1\cdots g_m)\le\sum_{j=1}^m d(\one,g_j),\qquad
 d(g^{-1},h^{-1})=d(g,h).
\]
If the eigenvalues of $g$ are $e^{\pm i\theta}$, $0\le\theta\le\pi$, then
\begin{equation}
 d(\one,g)=2\sin(\theta/2),\qquad
 s(g)=\tfrac12 d(\one,g)^2.\label{eq:energy-metric}
\end{equation}
Let $C_\rho$ be the geodesic cap $\theta\le\rho$, for $0<\rho\le1$.
Every $g\in C_\rho$ obeys $d(\one,g)\le\rho$ and is $e^X$ for a
traceless anti-Hermitian $X$ with $\norm X_{\mathrm{op}}\le\rho$.
Probability Haar measure is uniform measure on the unit three-sphere, so
\begin{equation}
 b(\rho):=\operatorname{Haar}(C_\rho)
 =\frac2\pi\int_0^\rho\sin^2\theta\dd\theta
 =\frac{\rho-\sin\rho\cos\rho}{\pi}
 \ge\frac{\rho^3}{24}.\label{eq:cap}
\end{equation}
The lower bound follows from $\sin\theta\ge\theta/2$ for
$0\le\theta\le1$ and $\pi<4$. These deliberately loose constants suffice.

\section{Why arbitrary frozen boundaries cannot be filled smoothly}
\begin{proposition}[A robust one-link obstruction]\label{prop:obstruction}
There is a frozen exterior for one editable link in the lattice
\eqref{eq:model} such that $E_J(U,z)=6$ for every $U\in\SU$.
For every $U$, at least three affected plaquettes obey
$d(\one,U_p)\ge\sqrt2$. Thus no replacement preserving that exterior can
make all affected plaquettes have curvature less than any $t<\sqrt2$.
The impossibility persists on an open set of frozen exterior data.
\end{proposition}
\begin{proof}
For one oriented link from $x$ to $y$, the six adjacent plaquettes can,
using trace cyclicity and inversion, be written with energies
$1-\tfrac12\Tr(U S_j^{-1})$, where $S_j$ is transport from $x$ to $y$
along the other three edges of that plaquette. These six staple paths
are internally edge-disjoint when the periods are at least four.
Set three staple transports to $+\one$ and three to $-\one$, assigning
central signs on a distinct edge of each path and identity on the rest.
These are actual link assignments, so plaquette constraints hold exactly.
Other, unaffected plaquettes may be rough; their contribution is constant
in the conditional law of the edited link.

Write $a=\Tr U/2\in[-1,1]$. The six energies are three copies each of
$1-a$ and $1+a$, with sum six. If $a\ge0$, the latter three are at
least one; if $a\le0$, the former three are. Apply
\eqref{eq:energy-metric}. The conditional law of $U$ is exactly Haar,
independently of $H$.

If each staple remains within chordal distance $\delta$ of its assigned
sign, bi-invariance and the reverse triangle inequality show that for
every $U$ the same three large-curvature plaquettes have curvature at
least $\sqrt2-\delta$. For fixed $t<\sqrt2$ choose
$0<\delta<\sqrt2-t$. Staple transports depend continuously on frozen
links, so this condition contains an open neighborhood of the displayed
boundary. Every nonempty open cylinder of link data has positive Haar
measure, and positive probability at every finite $H$ under the finite
Wilson law. No lower bound on that probability uniform in $H$ is asserted.
\end{proof}

\begin{proposition}[An explicit one-link boundary-defect quantity]
Identify $SU(2)$ with unit quaternions. For the six frozen staple
transporters let $s_j\in S^3\subset\R^4$, let $B=\sum_{j=1}^6s_j$, and
put
\[
 \mathcal D(z)=\sum_{i<j}d(S_i,S_j)^2.
\]
Then the boundary-forced minimum is exactly
\begin{equation}
 e_*(z)=6-|B|=\frac{\mathcal D(z)}{6+|B|},\qquad
 \frac{\mathcal D(z)}{12}\le e_*(z)\le\frac{\mathcal D(z)}6.
 \label{eq:staple-defect}
\end{equation}
The quantity $\mathcal D$ is gauge invariant and consists of frozen-loop
holonomy mismatches, since
$d(S_i,S_j)=d(\one,S_iS_j^{-1})$.
\end{proposition}
\begin{proof}
The fundamental normalized trace of $US_j^{-1}$ is the Euclidean inner
product $u\cdot s_j$, so $E_J(U,z)=6-u\cdot B$ and minimization over
$|u|=1$ gives $6-|B|$. Chordal operator distance of two quaternion
matrices is their Euclidean quaternion distance. Expanding the pairwise
squares gives $\mathcal D=36-|B|^2$, proving the identity and the bounds
using $0\le|B|\le6$. Gauge transformations act on all staples by the same
left and right unit quaternions, hence orthogonally on $\R^4$.
\end{proof}
This identifies a boundary quantity to retain or charge in the simplest
case. It does not show that it is statistically rare or decompose the
multi-link minimum into independent one-link quantities. Minimizing
aggregate energy also need not minimize the largest individual curvature.

The obstruction is boundary frustration, not an appeal to continuum
instanton topology. An unrestricted finite lattice has a connected
product of $SU(2)$ link spaces. A gauge transformation given at finitely
many boundary vertices can always be assigned arbitrary values at the
remaining vertices. A nonzero-degree continuous map $S^3\to SU(2)$
therefore does \emph{not}, by itself, prohibit a discrete link filling.
Continuum extension obstructions become relevant only with additional
regularity, interpolation, admissibility, or sector constraints, which
must be specified separately.

\section{An arbitrary-boundary excess-energy estimate}
Set $e_*(z)=\min_x E_J(x,z)$; the minimum exists by compactness.
It is continuous in $z$, since the uniform difference of the two energy
functions bounds the difference of their minima.

\begin{lemma}[A cap near a conditional minimizer]\label{lem:minimizer}
For every frozen $z$, choose any minimizer $V$ of $E_J(\cdot,z)$.
For $x_e=V_e e^{X_e}$ with $\norm{X_e}_{\mathrm{op}}\le\rho$,
\[
 E_J(x,z)\le e_*(z)+12\sum_{e\in J}\norm{X_e}_{\mathrm{op}}^2
 \le e_*(z)+12n\rho^2.
\]
\end{lemma}
\begin{proof}
Take the simultaneous path $x_e(t)=V_e e^{tX_e}$. A reversed occurrence
has inverse $e^{-tX_e}V_e^{-1}$. Differentiating a plaquette product twice
and using unitary operator norms gives
\[
 \left|\frac{\dd^2}{\dd t^2}s(U_p(t))\right|
 \le\left(\sum_{e\in\partial p\cap J}\norm{X_e}_{\mathrm{op}}\right)^2.
\]
Here $|\Tr A|/2\le\norm A_{\mathrm{op}}$, including nonsymmetric
complex matrices. Cauchy--Schwarz and the incidence bound imply
\[
 |E_J''(t)|\le\sum_p q_p\sum_{e\in\partial p\cap J}\norm{X_e}^2
 \le24\sum_{e\in J}\norm{X_e}^2.
\]
The first derivative at zero is zero because $V$ minimizes on the smooth
compact product manifold. Integrate the second derivative with weight
$1-t$ from zero to one.
\end{proof}

\begin{theorem}[Uniform excess above the boundary minimum]\label{thm:excess}
For every frozen exterior $z$ and $h\ge0$,
\begin{equation}
 \mu_H\{E_J\ge e_*(z)+h\mid z\}
 \le \min\{1,\exp[-Hh+n(12+\log24+\tfrac32\log H)]\}.
 \label{eq:excess}
\end{equation}
\end{theorem}
\begin{proof}
At $\rho=H^{-1/2}$, the product cap in Lemma~\ref{lem:minimizer} has
Haar mass $b(\rho)^n$ and energy at most $e_*(z)+12n/H$. It gives the
conditional partition-function lower bound
$b(\rho)^n e^{-He_*(z)-12n}$. The event numerator is at most
$e^{-H(e_*(z)+h)}$, since its Haar mass is at most one.
Divide and apply \eqref{eq:cap}. This pointwise argument uses no
measurable choice of a minimizer as $z$ varies.
\end{proof}
This theorem does not bound $e_*(z)$ itself. In
Proposition~\ref{prop:obstruction} it equals six and every link value is
already a minimizer. The theorem is therefore compatible with rough
plaquettes having conditional probability one.

\section{A sufficient boundary compatibility condition}
Let $F_J$ be the set of nonedited links occurring on plaquettes in $P(J)$,
and let $W_J$ contain all endpoints of links on those plaquettes. Form the
frozen graph $(W_J,F_J)$, including isolated vertices. Choose a spanning
forest, one root per component, and path orientations by a deterministic
ordering depending only on $J$ and lattice labels, never on link values.
For vertex $x$ let $T_x(z)$ be ordered transport from its component root
to $x$ along the tree, and set $h_x=T_x^{-1}$. For each frozen link
$e=(x,y)$ define
\begin{equation}
 R_e(z)=h_x^{-1}z_eh_y,\qquad
 \eta_J(z)=\max_{e\in F_J}d(\one,R_e(z)),\label{eq:residual}
\end{equation}
with the maximum over an empty set defined as zero. Tree residuals are
identity. All other residuals are fundamental-cycle holonomies.

\begin{lemma}[Meaning and covariance of compatibility]\label{lem:compatibility}
The function $\eta_J$ is continuous and gauge invariant. The condition
$\eta_J=0$ says exactly that every frozen-graph cycle has identity
holonomy, equivalently that there are vertex frames satisfying
$z_{xy}=h_xh_y^{-1}$ on all frozen links. For every $\eta>0$, the event
$\eta_J(z)<\eta$ contains a nonempty open neighborhood of identity
frozen data.
\end{lemma}
\begin{proof}
For site gauges $g_x$, path transports transform as
$T_x(z^g)=g_{r(x)}T_x(z)g_x^{-1}$. Thus
$h_x(z^g)=g_xh_x(z)g_{r(x)}^{-1}$ and every frozen residual transforms
by conjugation at its component root. Continuity and metric invariance
prove the first claim. Tree paths and fundamental cycles prove the flat
claim directly. At identity data all residuals are identity; continuity
proves the last assertion.
\end{proof}
The condition depends on the chosen forest quantitatively, not through
gauge choice. It is sufficient, not necessary, for low-curvature filling.
A long loop can have nontrivial holonomy while each individual plaquette
in a spanning surface is small. Exact flat data may be a measure-zero
class; the positive-tolerance version below avoids relying on that class.
No estimate on how often a useful tolerance is satisfied is supplied.

\section{A gauge-covariant link replacement kernel}
\begin{theorem}[Compatible-boundary cap replacement]\label{thm:kernel}
For $0<\rho\le1$ and fixed $J$, there is a Borel replacement kernel
$K_z(x,\dd y)$ with these properties:
\begin{enumerate}
\item All links outside $J$ stay exactly fixed. The conditional law of the
new edited links $y_J$ depends only on the frozen links in $F_J$, not
on the original $x$. The full-output law additionally copies the exterior.
\item Its conditional density in $y_J$ with respect to product Haar is at
most $b(\rho)^{-n}\le(24\rho^{-3})^n$.
\item Under every site gauge transformation it pushes forward to the
kernel for the transformed data.
\item For every output and each $p\in P(J)$,
\begin{equation}
 d(\one,y_p)\le q_p\rho+(4-q_p)\eta_J(z).\label{eq:output-curvature}
\end{equation}
If $\eta_J(z)\le\rho$, then every affected curvature is at most $4\rho$
and $E_J(y,z)\le48n\rho^2$. If $\eta_J(z)=0$, the sharper bound
$E_J(y,z)\le12n\rho^2$ holds.
\end{enumerate}
\end{theorem}
\begin{proof}
For each connected component $C$ of the frozen graph choose an independent
Haar frame $A_C$. Put $S_x=h_xA_{C(x)}$.
Independently choose each $Q_e$, $e\in J$, with normalized Haar restricted
to $C_\rho$, and set
\begin{equation}
 y_e=S_xQ_eS_y^{-1}\quad(e=(x,y)\in J),\qquad
 y_e=z_e\quad(e\notin J).\label{eq:kernel}
\end{equation}
All maps are finite products of continuous maps and Borel cap indicators,
so the law is a Borel kernel. Given the component frames, left/right Haar
invariance gives joint density at most $b(\rho)^{-n}$; averaging frames
preserves that bound. There is no division by a gauge-orbit volume.

Under a gauge transform couple $A_C^g=g_{r(C)}A_C$ and keep the same
$Q_e$. The transformed frames satisfy $S_x^g=g_xS_x$, while the coupled
root frames retain their independent Haar laws. Thus
$y_{xy}^g=g_xy_{xy}g_y^{-1}$, proving covariance.

In the $S$ frames, edited links equal $Q_e$ and each frozen link equals
$A_C^{-1}R_eA_C$. The metric triangle inequality around a plaquette gives
\eqref{eq:output-curvature}, including reversed edge orientations.
If $\eta_J\le\rho$, use
$s(y_p)\le(4\rho)^2/2$ and $|P(J)|\le6n$. For exact flat data, use
$s(y_p)\le q_p^2\rho^2/2$ and \eqref{eq:incidence}.
\end{proof}
Since the outputs are link configurations, all plaquette identities hold
automatically. Every loop consisting entirely of unedited links is
preserved exactly. No claim is made about preservation of loops traversing
$J$, an instanton number defined on an admissible interpolation, or any
sector label not included in the frozen data. No reflection-positivity
claim is made for the measure pushed forward by this kernel.

\section{The integrated energy--entropy estimate}
The following elementary measure argument is the useful replacement for
an unjustified invertible ``surgery.'' It tracks lost information by a
density bound, not by asserting a change-of-variables inverse.

\begin{lemma}[Bounded-density replacement]\label{lem:integrated}
Let $x=(x_J,z)$ have product probability reference measure. Let
$\mathcal Q\ge0$ be measurable with
$Z_{\mathcal Q}=\int e^{-HE(x)}\mathcal Q(x)\dd x\in(0,\infty)$.
Suppose a probability kernel changes only $J$, has conditional density
at most $M$, and for all original $x$ in a measurable event $\mathcal A$
and almost all outputs $y$ obeys
$E(x)-E(y)\ge h$ for one fixed real $h$. Then
\begin{equation}
 \frac1{Z_{\mathcal Q}}
 \int_{\mathcal A}e^{-HE(x)}
       \E_{K_x}[\mathcal Q(y)]\dd x\le M e^{-Hh}.
 \label{eq:integrated}
\end{equation}
For $\mathcal Q=1$ the left side is $\mu_H(\mathcal A)$.
For the probability under the $\mathcal Q$-weighted law, additionally
assuming $\mathcal Q(x)\le e^{Ln}\mathcal Q(y)$ on these pairs yields
$\mu_{H,\mathcal Q}(\mathcal A)\le e^{Ln}M e^{-Hh}$.
\end{lemma}
\begin{proof}
On the kernel support replace $e^{-HE(x)}$ by
$e^{-Hh}e^{-HE(y)}$. Bound the kernel density by $M$ and then discard
all restrictions on the original $x_J$ and the frozen $z$. Tonelli's
theorem applies to the nonnegative integrand. Integrating the original
$x_J$ contributes at most one; the remaining integral in $(y_J,z)$ is
$Z_{\mathcal Q}$. For the last statement use the extra comparison of
$\mathcal Q(x)$ and $\mathcal Q(y)$ before applying the same argument.
\end{proof}
No probability claim for an arbitrary $\mathcal Q$ follows by silently
identifying $\mathcal Q(x)$ with $\mathcal Q(y)$. Signed or complex
perturbations require an appropriate positive base weight and a separate
absolute-value/control argument. This lemma does not supply the
renormalized model's interaction bounds or its covering expansion.

Let $p\ge4$, let $k$ be a positive integer, and set
$\rho=H^{-1/2}$, $t=p/\sqrt H$.
Define the event, for fixed $J$,
\[
 \mathcal A_{J,k}=
 \{\eta_J(z)\le H^{-1/2},\quad
   |\{p'\in P(J):d(\one,U_{p'})>t\}|\ge k\}.
\]
When $t\ge2$ the rough-plaquette event is empty; the interesting regime
is $p<2\sqrt H$.

\begin{theorem}[Compatible rough-plaquette suppression]\label{thm:bad}
For each boundary with $\eta_J(z)\le H^{-1/2}$, the conditional probability
of at least $k$ rough affected plaquettes has the following bound. The
same bound holds for the joint event $\mathcal A_{J,k}$ under the finite
Wilson law:
\begin{equation}
 \mu_H(\mathcal A_{J,k})
 \le\min\{1,\exp[-kp^2/2+n c_H]\},\qquad
 c_H=48+\log24+\tfrac32\log H.\label{eq:bad}
\end{equation}
The exponential bound, without truncation by one, holds for the
replacement integral \eqref{eq:integrated} with any nonnegative
$\mathcal Q$ and this event. That integral need not be a probability.
All originally rough affected plaquettes are removed by the output, and
no affected plaquette exceeds $t$. Unaffected plaquettes do not change.
For exactly flat compatible boundary data, replace 48 by 12 in the
conditional bound.
\end{theorem}
\begin{proof}
On the event, \eqref{eq:energy-metric} gives
$HE_J(x,z)\ge kp^2/2$. Theorem~\ref{thm:kernel} gives
$HE_J(y,z)\le48n$ and density at most $(24H^{3/2})^n$.
Unaffected energies cancel exactly. Apply Lemma~\ref{lem:integrated}.
The output curvature is at most $4/\sqrt H\le t$.
For a fixed boundary the same proof uses only product Haar on $J$.
\end{proof}
The signed/nonlocal factors generated by renormalization are not assumed
to meet this theorem's raw-Wilson energy bound. The displayed general
$\mathcal Q$ integral is a restricted interface, not a proof that they do.

\section{Edited volume, counting, and what is still missing}
One bad plaquette does not imply original energy proportional to the
number of edited links. For example, in identity background set one link
to $-\one$ and take $J$ to be an arbitrarily large connected link set
containing it. Only six plaquettes are nonidentity, and $E_J=12$
independently of $|J|$. Thus the density cost cannot be charged to
edited volume without a separate geometric energy-density statement.

\begin{corollary}[A scoped animal bound]\label{cor:animals}
Suppose an event for each fixed connected $J$ imposes its compatibility
condition and at least $|J|/s$ rough affected plaquettes, where $s\ge1$.
If $p^2\ge4s c_H$, its probability is at most
$\exp[-|J|p^2/(4s)]$.
Connect two lattice links when they share a plaquette. For a fixed root
link and $\kappa=p^2/(4s)>2\log18$, the probability of at least one such
connected event containing that link is at most
\[
 \frac{e^{-\kappa}}{1-18^2e^{-\kappa}}.
\]
\end{corollary}
\begin{proof}
Apply \eqref{eq:bad} with $k\ge n/s$ and absorb $nc_H$ using the
assumed budget. The link adjacency graph has degree at most $6\cdot3=18$.
Every connected set of $n$ links containing the root has a deterministically
chosen spanning-tree traversal of length $2(n-1)$; the traversal determines
the visited set. Thus there are at most $18^{2(n-1)}$ such sets.
The union bound and a geometric series prove the claim; no independence
of the region events is used.
\end{proof}
For $p=(\log H)^P$ and $s=(\log H)^q$, the elementary absorption is
available for large $H$ if $2P>q+1$. This numerical budget is meaningful
only after the asserted compatible regions and bad-plaquette density
have actually been obtained.

Disjoint edited link sets do not suffice for additive local estimates
when their affected plaquettes overlap. If the sets $P(J)$ are disjoint,
energies and product-kernel density costs can be combined; otherwise one
must use the union with its actual affected plaquettes. For regions
chosen from the original configuration, freeze their geometry in each
integration term and sum possible geometries explicitly. The animal
bound covers the listed events, not an unconstructed decomposition of
all rough configurations.

\subsection*{Unresolved inputs for the cutoff-comparison program}
\begin{enumerate}
\item A local region-selection or defect-label scheme must handle typical
\emph{incompatible} boundaries. The frustrated example cannot simply be
excluded; its probability or alternative retained energy must be bounded.
\item Such a scheme must supply enough original energy per edited link
and controlled counts of regions and boundary labels. The raw one-link
estimate and the forest condition do not do this automatically.
\item Exact renormalization generates interactions beyond the original
Wilson sum. Their changes under replacement, locality, and signed
covering weights need independent uniform estimates.
\item Only after those estimates could one seek a cutoff-comparison
error small on a fixed physical box. Continuum observables and their
reconstruction remain additional tasks, as in the first research note.
\end{enumerate}
The present results settle the proposed \emph{unrestricted local filling}
in the negative and give two precise restricted replacements. They do
not establish the nonperturbative cutoff-comparison lemma or resolve
Yang--Mills existence and mass gap.

\section{Verification scope}
The accompanying audits independently check the boundary obstruction,
frame covariance and density, and integrated estimates. The Python
checks exercise the link geometry, matrix identities, gauge covariance,
curvature bounds, incidence constants, and frustrated negative controls.
Numerical tests are not proofs of the theorems. Any Lean certificate is
restricted to the stated scalar arithmetic; no analytic, gauge-theory,
or continuum formalization is claimed. The first research pass and
published repository sources are preserved. Literature context and
precise source-access limits are in \texttt{audits/literature.md}.

\begin{thebibliography}{9}
\bibitem{repo-o3}
OpenAI, \emph{The canonical massive continuum limit of the two-dimensional
$O(3)$ model}, October 4, 2026. Supplied repository,
\texttt{build/sections/comparison.tex}, Lemma ``Replacement and its energy
bounds'' and equation labeled \texttt{eq:orig-22}.
\bibitem{wilson}
K. G. Wilson, \emph{Confinement of quarks}, Physical Review D
\textbf{10} (1974), 2445--2459.
\url{https://doi.org/10.1103/PhysRevD.10.2445}.
\bibitem{clay}
A. Jaffe and E. Witten, \emph{Quantum Yang--Mills Theory}, official Clay
problem description.
\url{https://www.claymath.org/wp-content/uploads/2022/06/yangmills.pdf}.
\end{thebibliography}
\end{document}
